\clearpage
\fancyhead[L]{\small BODE. ASTRONOMICAL YEARBOOK.}
\hypertarget{native-624-lower}{}
\begin{center}\small Göttingische gelehrte Anzeigen. No. 23. Pages 225..232. 7 February 1820.\\[3pt]\rule{.72\linewidth}{.3pt}\end{center}

\emph{Astronomical Yearbook for the Year {\upshape 1822}, together with a collection of the latest papers, observations, and reports pertaining to the astronomical sciences. Calculated and published by} Dr. J. E. \textsc{Bode}, \emph{Royal Astronomer, etc.} Berlin. Available from the author and on commission through F. Dümmler. 260 pages in octavo, with one copperplate.

\begingroup\rightskip=0pt plus 1em
Of the two solar and two lunar eclipses that occur in 1822, only one of the lunar eclipses is visible in Europe. The transit of Mercury on 5 November likewise cannot be observed in Europe. The arrangement of the yearbook has remained the same as in the preceding volumes. The only change is that, to make the work easier, the eclipses of Jupiter's satellites have been calculated using \textsc{Wargentin}'s tables, with \textsc{Delambre}'s tables occasionally used as a supplement for the first satellites alone.\par\endgroup

\begingroup\brokenpenalty=10000\relax
The first paper in the appendix is an attempt by Prof. \textsc{Buttmann} to establish the nomenclature of individual stars. This learned philologist \SourceLeaf{625}presents here a list of 80 star names, mostly borrowed from Arabic. Some were already in use or at least were not unknown, others had been corrected or reassigned, and still others had been newly coined (e.g., Clavus for $\beta$ in Taurus, Meion for $\varkappa$ in Orion, 
Sectatrix for $\beta$ in Auriga). Among astronomers, only a very small number of proper names for the fixed stars are \emph{actually in use} (about a dozen), and attempts by some to reintroduce older Arabic names have found no favor. Should one nevertheless wish to bring a greater number into use, the most practical course would probably be simply to retain \emph{without any alteration} the names that \textsc{Piazzi} included in his new star catalogue. Any arbitrary change could lead only to confusion in the nomenclature instead of improving it. --- Eclipses of Jupiter's satellites, occultations of stars, lunar eclipses, and oppositions of Vesta, Uranus, Jupiter, Saturn, and Pallas, observed in 1818 by Prof. \textsc{Bürg}. --- Algol's period of light variation determined from more recent observations, and calculation of its minimum light for 1820 to 1822, by Prof. \textsc{Wurm}. This study shows that the length of the period has undergone no perceptible change at all since the discovery of the star's variability. --- Astronomical observations at the Prague Observatory by Messrs. \textsc{David} and \textsc{Bittner}. --- Also by Prof. \textsc{Bürg}, observations of Juno and occultations of stars in 1819, together with astronomical reports. The defense presented here against an attack on or a misinterpretation of a newspaper article concerning one of the comets of 1818 (the attack or misinterpretation appeared in the Correspondance astronomique published in Genoa) prompts the wish that personal remarks may always be kept out of astronomical discussions. --- Occultations of stars and eclipses of Jupiter's satellites observed in Prague by Prof. \textsc{Hallaschka}. --- Earlier observations of Uranus by \textsc{Lemonnier} (see the 1819 volume, p. 1909). --- On the solar eclipse of 7 September 1820 for Germany and the adjoining countries, by Prof. \textsc{Littrow}. The beginning and end of the eclipse are given for 230 places in Germany and neighboring countries, as are the northern and southern boundary lines, defined by a number of points, of the zone within which the eclipse will appear annular. The greater part of Germany lies within this zone. The southern boundary passes near Aachen, Strasbourg, Freiburg, and Zurich, and the northern boundary near Schwerin, Wittenberg, and Tabor in Bohemia. It will be very desirable for the duration of the annular eclipse to be observed at a good many places, especially those near the boundaries, .for which all that is needed is a good telescope and a seconds clock. --- Observations of Vesta, Pallas, Uranus, Saturn, and Jupiter, and of the solar eclipse in 1818, at Wilna by Prof. \textsc{Sniadecky}. --- On the geographical longitude of the Berlin Observatory. This longitude was formerly always counted among those whose determinations varied widely. Since the observatory was equipped with better means of determining the time, however, the observations too have shown more satisfactory agreement. From 23 stellar occultations and 4 solar eclipses observed from 1802 to 1815, there follows a mean longitude of $44^{\mathrm{m}}10^{\mathrm{s}}5$ from Paris. --- Observed oppositions of Mars, Uranus, and Jupiter in 1817, 1818, and observations of the solar eclipse of 1818 and of \textsc{Enke}'s comet, by Mr. \textsc{Derflinger} in Kremsmünster. --- Observation of an occultation of Mars by the Moon, made in Vienna by Prof. \textsc{Bürg}. --- Astronomical observations at the Berlin Observatory in 1818. The hasty assertion made here needs correcting. It claims that determining the heliocentric position from a geocentric position observed near opposition becomes uncertain because of the smallness of the angles of the triangle, and that it is even preferable to base the determination on observations farther from opposition. --- Observations of a comet in 1818 and 1819 by Prof. \textsc{Bessel}, as well as planetary oppositions, a solar eclipse, and solstices. --- By Dr. \textsc{Olbers}, on the highly remarkable discovery \SourceLeaf{626} made by Prof. \textsc{Enke} and first announced in our pages (St. 28 u. 83 v. v. J. [S. 417 u.f. d. B.]), that one of the comets of 1805 was identical with one of the comets of 1818. Dr. \textsc{Olbers} added to this great discovery the equally important observation that the comet of 1795 was also very likely the same comet, and that the two comet observations published in the Connaissance des tems 1819 and dating from 1786 likewise pertain to that very comet. This distinguished astronomer also reports here his first observations of the comet that appeared in July of this year. --- The fine article that follows, by Prof. \textsc{Enke}, gives a detailed account of his work on that comet with a short orbital period and raises the probable conjecture by Dr. \textsc{Olbers} mentioned above to moral certainty. In mid-May 1822 this comet will again reach perihelion. Unfortunately, because of its position relative to the Earth, it will be extremely faint beforehand, so it will be difficult to observe it until then. After its passage through perihelion, on the other hand, it will be bright enough, but it will be too far south to become visible in Europe. It will then be easy to observe from the southern half of the globe, if only there are observers there. On occasions like this one, as on so many others, the need for a \emph{southern} observatory is keenly felt. On this occasion we feel compelled to express our deep conviction that, after the great support astronomy has received in Europe in recent times, no greater service could now be rendered to this science than by establishing a well-equipped observatory in the southern hemisphere, perhaps at the Cape of Good Hope, to be staffed by a perceptive, experienced, and active astronomer. --- By the same astronomer, observations of the comet of 1819 and elements of its orbit, as well as observations of Uranus, Saturn, Jupiter, Ceres, and Juno in 1818 and 1819. --- Solution by Prof. \textsc{Fischer} to the problem of finding the declination of a celestial body with a telescope alone, without an instrument for measuring angles (from the duration of a transit through the field of view in which entry and exit occur at the endpoints of a diameter). --- Geometrical motion of Ceres in 1819 and 1820, calculated by Dr. \textsc{Westphal}, and geocentric motion of Juno in 1820, calculated by Prof. \textsc{Nicolai}. --- Observations of the comet of July 1819 at the Berlin Observatory (running from 2 to 27 July). --- By Prof. \textsc{Nicolai}, observations of Juno at the Mannheim Observatory, and new orbital elements derived from these and from observations at Göttingen. Also, the first observations of the Pole Star with the restored three-footed \textsc{Reichenbach} repeating circle. These are in very close agreement and give the altitude of the pole at the Mannheim Observatory as $49^\circ\,29'\,12''95$. Finally, observations and elements of the comet of 1819. --- The ensuing remarks by an unnamed author on using the spirit level and the plumb line to adjust astronomical instruments contain nothing new and show only that the way spirit levels are mounted on an instrument matters, that their sensitivity must suit their purpose, and that they require a careful observer. The proposal made here to use a plumb line suspended in a well to test whether the direction of gravity varies or remains constant would scarcely do more than demonstrate the \emph{impossibility of attaining absolute stability} in any setup of measuring apparatus. --- Extract from a letter from General von \textsc{Lindener} of Glatz. By a fortunate coincidence, the inferior conjunction of the comet that appeared in July of this year coincided with its passage through the ascending node on the morning of 26 June. It was therefore most welcome that this keen enthusiast for astronomy had indeed observed the Sun at that very time and could thus, by actual experience, \SourceLeaf{627} confirm what we could only have expected, given everything else we presently know, namely that, although the comet's diffuse body was certainly and exactly between the Earth and the Sun, not even the slightest trace of it was discerned on the solar disc. We recall reading in the newspapers that another observer in Vienna, who was also observing the Sun during the same hour, noticed nothing whatsoever out of the ordinary. --- Observations of this comet in the latter half of August by Dr. \textsc{Olbers}, as well as elements of its orbit by Mr. \textsc{Bouvard} and Mr. \textsc{Dirksen}. --- On the problem of finding the position of a celestial body from observed alignments, by Dr. \textsc{Olbers}. The Yearbook for 1821 contains a solution to this problem by \textsc{Bessel}. The form chosen by \textsc{Olbers} requires two fewer logarithms. It should be remembered, however, that the 20 logarithms required for this form are found at as many different locations, whereas the 22 needed for \textsc{Bessel}'s form need to be looked up at only 17 different locations. We also note that if the familiar auxiliary table for logarithms of sums and differences is used as well, one can manage with 18 logarithms at 16 locations, or with 20 logarithms at 14 locations. --- Observations of the comet of July 1819, a description of \textsc{Reichenbach}'s six-foot meridian telescope, and an observation of Jupiter's opposition and of the Pole Star, by Court Councillor \textsc{Gauss}. In connection with what is said here about the absolute reliability of the zenith distances, we draw attention to two misprints that completely distort the meaning. At the top of p. 239 [p. 427], for merklich read \emph{unmerklich}, and in the 5th line from the bottom, for doch kann ich schon glauben read \emph{doch kann ich schwer glauben}. --- Observations of the comet of July 1819 at the Kraków Observatory by Prof. \textsc{Leski}. --- List of the precise longitude and latitude of nine of the principal fixed stars for which lunar distances appear in the Nautical Almanac and in the Connoissance des tems, taken from the latter's volume for 1821 (see last year's Anz., p. 1907 [p. 615 of this volume]). --- Observations of a new comet made in Marseille, and elements of its orbit calculated by Prof. \textsc{Enke}. So far as is known, this comet was observed only in Marseille in June 1819, and it is quite remarkable that the observations can be represented much better by an ellipse with an orbital period of $2\frac{1}{2}$
years than by a parabola. --- Observations of the (other) comet of 1819 at the Greenwich Observatory, and parabolic elements derived from them by Mr. \textsc{Rumker}. --- The brief reports that follow at the end contain several more astronomical observations and notes, which cannot be listed individually here. The accompanying copperplate depicts the geocentric course of the comet of July 1819 and a projection of its true orbit, as well as projections of the orbits of \textsc{Enke}'s comet and of the other comet observed from November 1818 to January 1819.
\par\clearpage\endgroup
